\subsection{Local Linear Convergence with Adaptive Restart}
\label{app:pduot-restart}

\begin{proof}[Local linear convergence in Proposition~\ref{prop:pduot-convergence}]
Global convergence places all sufficiently late iterates in the KKT error-bound neighborhood of $\mathbf z^\infty$ from~\eqref{eq:pduot-kkt-error-bound}.
The uniform spectral bounds transfer that error bound to every late-epoch map, giving a common constant $\alpha_0>0$ such that
\begin{equation}
\label{eq:pduot-uniform-sharpness}
\alpha_0\operatorname{dist}_{\mathbf P_j}(\mathbf z,\mathcal S)\leq\|\mathbf z-T_j\mathbf z\|_{\mathbf P_j}
\end{equation}
near the limit.
One may take $\alpha_0=(1+\kappa\overline\lambda)^{-1}$ by the argument for~\eqref{eq:pduot-local-sharpness}; the common neighborhood follows from nonexpansiveness and uniform norm equivalence.
Write $d_j=\operatorname{dist}_{\mathbf P_j}(\mathbf a_j,\mathcal S)$ and $\mathcal R_j(\mathbf z)=\|\mathbf z-T_j\mathbf z\|_{\mathbf P_j}$.
Nonexpansiveness gives $\mathcal R_j(T_j\mathbf t_j)\leq e_j$.
Applying~\eqref{eq:pduot-uniform-sharpness} at $\mathbf a_{j+1}=T_j\mathbf t_j$, and taking a nearest solution to $\mathbf a_j$ in~\eqref{eq:pduot-epoch-descent}, yields
\begin{equation*}
\operatorname{dist}_{\mathbf P_j}(\mathbf a_{j+1},\mathcal S)^2+e_j^2\leq d_j^2,\qquad \alpha_0\operatorname{dist}_{\mathbf P_j}(\mathbf a_{j+1},\mathcal S)\leq e_j.
\end{equation*}
With $q_0=(1+\alpha_0^2)^{-1/2}<1$, metric comparison then gives
\begin{equation*}
d_{j+1}\leq e^{c_0b_{j+1}/2}q_0d_j.
\end{equation*}
Since $b_j\to0$, for any $q\in(q_0,1)$ the right-hand factor is at most $q$ in all sufficiently late epochs.
Moreover $\|\mathbf a_{j+1}-\mathbf a_j\|_2\leq2d_j/\sqrt{\underline\lambda}$ by the within-epoch distance bound.
Summing these geometrically decreasing increments proves R-linear convergence to the actual limit, uniformly for all iterates and mapped points within each late epoch, when indexed by epoch.

For the work bound, let the sufficient-reduction factor be $\vartheta\in(0,1)$ and let $h$ bound the gap between restart checks, including the first check of each epoch.
The anchor baseline is evaluated under the current metric.
The Halpern estimate~\eqref{eq:pduot-halpern-bound} and uniform sharpness give
\begin{equation*}
\mathcal R_j(\mathbf z^{j,k})\leq\frac{d_j}{k+1}\leq\frac{\mathcal R_j(\mathbf a_j)}{\alpha_0(k+1)}.
\end{equation*}
Thus the practical sufficient-reduction condition
\begin{equation*}
\mathcal R_j(\mathbf z^{j,k})\leq\vartheta\mathcal R_j(\mathbf a_j)
\end{equation*}
holds once $k+1\geq1/(\alpha_0\vartheta)$.
Every sufficiently late epoch therefore has at most $H=\lceil1/(\alpha_0\vartheta)\rceil+h$ primary PDHG evaluations; additional restart triggers can only shorten it.
A final infinite epoch from a nonsolution is impossible, because its residual tends to zero and eventually passes the sufficient-reduction test.
If an anchor is already a solution, every subsequent iterate equals it.
Otherwise, after $N$ primary evaluations in the tail at least $\lfloor N/H\rfloor$ epochs have completed, giving a bound $Cq^{\lfloor N/H\rfloor}$ for the iterates, mapped points, and projection multipliers.
This proves the stated local linear rate in work.
\end{proof}

The practical factors $\vartheta=0.2$, moderate reduction $0.5$, and artificial epoch-length threshold $0.36$ are covered by this argument.
Periodic checks with a fixed interval satisfy the bounded-gap requirement even when not aligned with epoch starts.
The results assume exact proximal evaluations and fixed base steps within each restricted call; finite stopping caps and floating-point errors retain their usual numerical scope.
